\section*{Órbitas en Schwarzschild \normalfont{(\text{órbitas en un plano}$\rightarrow \theta =\pi/2\; \; \; (x^0, x^1, x^3)\equiv (ct, r, \phi)$)}}
\begin{align*}
&\text{Defino: }\breve{r}=\frac{r}{a}\; \; \; \breve{t}=\frac{ct}{a}\; \; \; \breve{\tau}=\frac{c\tau}{a}\; \; \; \breve{E}=\frac{E}{mc^2}\; \; \; \breve{L}=\frac{L}{mac};\; \;\; c=1 \; \; \; a=1\\
&\text{A partir de ahora, no usaré el símbolo $\smallsmile$  para denotar las unidades adimensionales.}\\
&\text{Para obtener los valores en S.I. debo aplicar las definiciones previas.}\\
&ds^2 = -\left(1-\frac{1}{r}\right)dt^2 + \frac{1}{1-\frac{1}{r}}dr^2 + r^2 d\phi^2=-d\tau^2 \; \;\;  \text{Killings: }\underline{\xi}=e_0\; \; \underline{\xi}=e_3\\
&\frac{dt}{d\tau}= \frac{E}{(1-\frac{1}{r})}\; \; \; \; \frac{d\phi}{d\tau}= \frac{L}{r^2}\;\; \; \Rightarrow \; \; \; \varepsilon = \left(\frac{dr}{d\tau}\right)^2 + \frac{L^2}{r^2 }- \frac{1}{r}- \frac{L^2}{r^3}\; \; \; \varepsilon\equiv E^2 -1\\
&\text{Ecuación diferencial de movimiento (S.I.): }\frac{d^2 u}{d\phi^2 }+ u = \frac{GMm^2 }{L^2}+ \frac{3GM}{c^2}u^2\; \; \; u\equiv\frac{1}{r}\\
&\text{Potencial efectivo: }V_{\text{eff}}(r)= \frac{L^2}{r^2}-\frac{1}{r}-\frac{L^2}{r^3}\; \;\;\;  \; 
\left|
\begin{array}{l}
r_{\text{c}1} = L^2+L\sqrt{L^2-3}\text{ (mínimo)}\\
 r_{\text{c}2} = L^2-L\sqrt{L^2-3} \text{ (máximo)}
\end{array}
\right.
\end{align*}
\begin{align*}
    &\bullet L^2 =0: \text{Órbitas radiales (movimiento en la línea que une los cuerpos).}\\
    &\bullet L^2=3: \text{Órbitas espirales}\\
    &\bullet 3<L^2<4: \text{Órbitas elípticas}
    \left|
\begin{array}{l}
V_{\text{mín}}<\varepsilon\leq V_{\text{máx}}: \text{Posible precesión}\vspace{4pt}\\
\varepsilon=V_{\text{mín}}: \text{Posible órbita circular}\\
\end{array}
\right.\\
    &\bullet L^2=4: \text{Órbitas elípticas}
        \hspace{13pt}
    \left|
\begin{array}{l}
V_{\text{mín}}<\varepsilon<V_{\text{máx}}: \text{Posible precesión}\vspace{4pt}\\
\varepsilon=V_{\text{mín}}: \text{Posible órbita circular}\\
\end{array}
\right.\\
    &\bullet L^2>4: \text{Órbitas elípticas}
    \hspace{13pt}
    \left|
\begin{array}{l}
V_{\text{mín}}<\varepsilon<0: \text{Posible precesión}\vspace{4pt}\\
\varepsilon=V_{\text{mín}}: \text{Posible órbita circular estable}\vspace{4pt}\\
\varepsilon=V_{\text{máx}}: \text{Posible órbita circular inestable}\\
\end{array}
\right.
\end{align*}
\begin{align*}
    &\text{Notas: solo podemos conocer la evolución del cuerpo hasta que atraviese el}\\
    &\text{horizonte de eventos ($r=1$), o colisione con otro cuerpo. $V(r\rightarrow\infty)=0$}
\end{align*}
